\chapter{记忆}

\section{记忆对象}

记忆是主体接受并决定长期保留的认知。每个记忆对象具有稳定身份，内容通过不可变修订演化；当前修订、对象生命周期和并发版本分别维护。

记忆形成可以由确定性流程或模型辅助流程提出。所有形成请求必须给出主张内容、语义指向、来源证据和形成上下文。模型生成的内容以提案进入记忆域，记忆域核验来源、当前对象状态和写入权限后提交。

\section{认知角色}

认知角色描述一条记忆在主体认知中承担的作用。它与证据数量和形成方式分开。

\begin{ReportWideTable}{>{\raggedright\arraybackslash}p{3.3cm}>{\raggedright\arraybackslash}p{10.1cm}}
\captionof{table}{记忆的认知角色}\label{tab:memory-roles}
\toprule
\textbf{认知角色} & \textbf{定义}\\
\midrule
经历性记忆（Experiential） & 保存主体经历、接触或参与过的事件，以及主体对该次经历的长期认知。\\
陈述性记忆（Declarative） & 保存主体长期接受的事实、命题、判断或解释。\\
程序经验（Procedural Experience） & 保存主体从经验中形成的做法、步骤或策略性认识。它属于认知，不直接构成可执行行为规则。\\
\bottomrule
\end{ReportWideTable}

\section{形成方式}

形成方式描述某一修订依据什么形成。

\begin{description}
\item[经历落地形成（Grounded）] 主要依据一个有界经历、观察事件或来源上下文形成，可以同时引用其他佐证。
\item[综合形成（Synthesized）] 依据多个认知对象、多个经历或多个独立来源综合形成。
\end{description}

认知角色和形成方式是独立属性。同一种认知角色可以采用不同形成方式；同一种形成方式也可以产生不同认知角色。对象身份不通过一个混合所有含义的单一类型枚举定义。

\section{修订与关系}

对同一认知的修正、澄清或重新解释创建新修订；独立的新事件或新认知创建新对象。当前修订表示该对象当前被采用的内容版本。

对象之间可以建立显式关系，包括推导、支持、矛盾、替代依据和其他具有领域语义的关系。矛盾关系不自动撤回任一对象；撤回表示领域所有者不再接受某项认知，必须由独立生命周期操作表达。

内容修订、对象生命周期、抑制状态和访问策略使用不同状态维度。仅改变抑制、固定访问或运行时驻留时，不创建虚假的内容修订。

\section{记忆可访问性}

记忆是否长期存在、当前是否允许普通访问、当前是否容易被召回以及与某次查询是否相关，分别由认知权威、管理生命周期、可访问性策略和查询相关性处理。

时间流逝本身不删除记忆。自动可访问性在查询时根据认知年龄、类型化使用、当前线索、工作上下文和查询预算计算；具体函数和参数由下游算法规格确定。精确管理读取和合法的来源追溯不因自动可访问性降低而丢失对象身份。
